\begin{center}
{\LARGE Maps from the complex six-spheres and the manifolds of the torus-fibration construction}\par\medskip
{\large Companion Paper III}\par\medskip
AI-produced by GPT-6 Astra (OpenAI, ChatGPT Pro mode, effort Pro),\par
17 November 2026
\end{center}
\begin{abstract}
For the standard marked complex six-sphere, the repaired circle quotient
has class $6\nu'$ in $\pi_6(S^3)=\mathbb Z/12$, the holomorphic map to
$\mathbb P^1$ is null-homotopic, and the positive parameter-loop smooth
monodromy is zero in $\Theta_7$. We compute the associated sphere maps,
framings, covers, cusp carriers, elliptic-surface constructions and fibre
products, retaining their integral boundary maps. A compact complex
fourfold is orientation-preservingly diffeomorphic to
$3\mathbb{HP}^2\#(-\mathbb{HP}^2)$. Complete per-input algorithms give
integral conjugacy representatives, fundamental groups and the homology
of marked rank-six gluings; their survivor criterion imposes no rational
height bound. Classical topology is credited and re-derived in full.
The independent proof readings, source checks and reruns were made by
GPT-6 Astra (OpenAI, ChatGPT Pro mode, effort Pro), 17 November 2026;
historical reviews retain their exact scopes. No human has checked the
mathematics. The remaining noncanonical six-dimensional filling and
refined-invariant questions are separated from the proved results.
\end{abstract}

\section*{Main statements, status, credit and proof map}
We use the family and marking of Paper I. The source sphere has its complex
orientation; normal frames precede tangential frames. The generator
$\nu'\in\pi_6(S^3)$ stabilizes to $2\nu$, where the real Adams invariant
of $\nu$ is $1/24$. The sphere-gluing orientation identifies the generator
of $\Theta_7$ with Eells--Kuiper invariant $1/28$.
The complete sphere-map statements are
\cref{maps:src:A.2.2,maps:src:B.1.2,maps:src:B.2.1,maps:src:B.3.1,maps:src:B.3.2,maps:src:B.3.3}.
The actual parameter-loop determination has the two complete arguments
\cref{push:src:C.3.24,push:src:C.3.27}.
The exact fourfold diffeomorphism is \cref{maps:src:F.4.1};
the effective integral classification with full marked gluings is
\cref{maps:src:F.2.39}.
These statements specify the actual maps, affine lifts, cusp model and
boundary markings on which their conclusions depend.

\subsection*{Status}
\begin{tabular}{@{}p{31mm}p{120mm}@{}}
\toprule
Kind & Theorems, computations, corrections, counterexamples and credited re-derivations.\\
Claim & Complete checked arguments in the body and appendices; exact unresolved targets in the separate final part.\\
Checks & Full independent readings and derivations, opened operative sources, and exact program reruns, at the individual scopes in the concordance.\\
Human understanding & None claimed; no human has checked the mathematics.\\
Formalization & None; the finite certificates are computational checks.\\
Open points & The remaining filling and refined-invariant cases in \cref{part:III:open}.\\
\bottomrule
\end{tabular}

\subsection*{Known and new}
\begin{tabular}{@{}p{42mm}p{70mm}p{37mm}@{}}
\toprule
Statement & Credited source and scope & Present treatment\\\midrule
The underlying complex six-sphere and elliptic-surface construction &
\citet[Theorem~10.5]{AC2026}; \citet[main theorem]{Engel2026} &
The family is fixed in Paper I; maps and relative constructions are specified here.\\
Low-dimensional framed and smooth classification inputs &
The exact Bott, Cerf, Eells--Kuiper, Wall, Nikulin and related statements cited at their uses &
Known conclusions credited; every distinct receiving derivation retained.\\
Integral lattice isomorphism algorithm &
\citet[Theorems~3.4, 7.9, 8.4; Algorithm~8.3]{rootBHJ22} &
Its hypotheses and the simultaneous-tuple reduction are proved explicitly.\\
Maps, marked gluings and manifold calculations &
The per-result source comparisons in Appendix~\ref{app:III:result-index} &
Exact proved scopes; no inference from a bounded sample to an infinite class.\\
\bottomrule
\end{tabular}

\subsection*{Proof map}
The primitive invariant covector gives the circle quotient and its framed
regular fibre. The full fibre and cap maps determine the unstable homotopy
classes. Two independent boundary and characteristic calculations determine
the parameter monodromy. Integral normalizations, exterior lattices and
actual attaching maps determine the higher-dimensional constructions.
The rank-five quotient normal form, finite conductor interval and primitive
extension calculation give all integral survivor candidates; local
suspension and torsion transgressions give the exact six-test criterion.
Marked chain models retain the full homology extensions for every input.
The appendices print the classical geometric, stable-homotopy and analytic
prerequisites and all distinct working-edition proofs. The concordance and
program index give exact statement and computation locations.

\clearpage
\section{Common conventions and the transport dictionary}\label{sec:III:conventions}
The base is $B=\mathbb P^1_t$, with $(p_1,p_2,p_0)=(0,1,\infty)$,
and the regular-fibre lattice is
$\Lambda=\mathbb Z\langle\widehat\gamma,\widehat u,\widehat w,\widehat\delta\rangle$.
The covector $\gamma$ is dual to $\widehat\gamma$.
The family homology matrices are $A_1,A_2,M_0$, with
$A_1A_2M_0=I$; their cohomology transport matrices are
$T_j=A_j^{-t}$. An action of $A_j$ on a cohomology class means its
induced pullback, not multiplication of a homology column by that class.
Base meridians are counterclockwise. The analytic root generators are
clockwise, $s_j(g_jz)=e^{-2\pi i/m_j}s_j(z)$, and
$g_1g_2g_0=1$. The retained multiplicities are $(3,4)$, with
$f^*p_j=m_jS_j$ and reduced cusp $W=f^{-1}(p_0)$.
We use
\[
\Pi_c=\begin{pmatrix}6\mu&\tau&1&0\\\beta^{\rm part}+c&\mu&0&1\end{pmatrix},
\qquad w=e^{-\pi i c},\qquad
D_{\rm part}=\operatorname{Im}\beta^{\rm part}
-\frac{6(\operatorname{Im}\mu)^2}{\operatorname{Im}\tau}.
\]
The normalization is $(\beta^{\rm part}+\tau)(p_0)=0$ and
$M=\max D_{\rm part}=-3\log(6)/\pi$.
For a fixed member $D=D_{\rm part}+\operatorname{Im}c$;
admissibility is $\operatorname{Im}c<-M$.
The positive $w$-loop lifts to $c(r)=c_0-2r$.

Engel's indices $(0,1,\infty)$ are our $(1,2,0)$.
His line bundle is denoted $M_S=\mathcal O_S(P-O)$.
The common marking matrix is
\[
P_E=\begin{pmatrix}0&-2&-1&-2\\2&-1&0&-1\\0&2&1&3\\-1&0&0&0\end{pmatrix},
\qquad \psi P_E=-\gamma,\qquad a_{j-1}=-P_Ev_j/m_j.
\]
The standard twists are $v_1=(1,2,-4,0)^t$ and
$v_2=(-1,-3,3,0)^t$. For general twists
$p=12l_0-4l_1-3l_2$ and $\pi_1(X_p)=\mathbb Z/|p|$.
The vector shorthand $(g,u,w,d)$ in fibre calculations denotes the ordered
homology basis, never the scalar parameter $w$ in a period expression.

\paragraph{Rank-six matrices.}
In the fourfold construction we use $M_i$ for actual homology transport
and $T_i=M_i^{-t}$ for left cohomology transport. The explicitly constructed
matrices named $A,B$ in the rational normal-form calculation are geometric
cohomology pullbacks: their homology inputs are $(A^t,B^t)$ and their left
cohomology transport is $(A^{-1},B^{-1})$. Every peripheral product is kept
in its stated order; reversing generators also conjugates or inverts the
cusp as specified in the proof. The fixed Kuhn cap and its full marked
boundary are part of a gluing datum. Affine multiplicity, linear order,
translation numerator and ordinary clutching are separate inputs.

All homology and cohomology have integral coefficients unless a coefficient
group is printed. We use the complex orientation, outward stabilization,
and the displayed ordered normal coordinates. Local applications of letters
for an integer matrix, a cap or a lattice are explicitly defined at their
use; the complete symbol and source-label dictionary accompanies the papers.
